\section{Od grupy do sfery bez punktów wyjątkowych}
Od tej chwili oznaczamy przez $G=\rho(F_2)\le\mathrm{SO}(3)$ wolną podgrupę z twierdzenia~\ref{tw:wolne-obroty}. Dzięki różnowartościowości możemy utożsamiać elementy $G$ z ich jedynymi słowami zredukowanymi. Litery $a,b$ w działaniu na sferze oznaczają obroty $A,B$.

\subsection{Orbity, stabilizatory i działanie swobodne}
\begin{definicja}
Orbita punktu $x$ to $\Orb(x)=Gx=\{gx:g\in G\}$. Stabilizator punktu to $\Stab(x)=\{g\in G:gx=x\}$. Działanie jest \emph{swobodne}, jeśli $\Stab(x)=\{e\}$ dla każdego punktu rozważanego zbioru.
\end{definicja}
Relacja „należą do tej samej orbity” jest równoważnością: $x=ex$, równanie $y=gx$ daje $x=g^{-1}y$, a z $y=gx$ i $z=hy$ mamy $z=(hg)x$. Orbity są więc parami rozłączne i pokrywają zbiór, na którym działa $G$.
\begin{uwaga}[Grupa wolna a działanie swobodne]
Wolna grupa nie musi działać swobodnie! Nasza grupa $G$ jest wolna, ale obrót $A$ ustala dwa punkty osi $x$ na sferze. Aby punkt sfery miał jednoznaczny opis przez słowo i reprezentanta orbity, trzeba usunąć wszystkie punkty ustalane przez element $G$ różny od identyczności.
\end{uwaga}

\subsection{Przeliczalny zbiór wyjątkowy}
Definiujemy
\begin{equation}\label{eq:6-1}
 D=\bigcup_{g\in G\setminus\{e\}}\Fix_{S^2}(g),\qquad X=S^2\setminus D,
\end{equation}
gdzie $\Fix_{S^2}(g)=\{x\in S^2:gx=x\}$.
\begin{lemat}\label{lem:wyjatkowy}
Zbiór $D$ jest co najwyżej przeliczalny i niezmienniczy względem $G$. Zbiór $X$ jest niezmienniczy, a działanie $G$ na $X$ jest swobodne.
\end{lemat}
\begin{proof}
Grupa $G$ jest przeliczalna na mocy stwierdzenia~\ref{st:przeliczalnosc-f2} i monomorfizmu $\rho$. Każdy $g\ne e$ jest nietrywialnym obrotem i ma dokładnie dwa punkty stałe na sferze, zgodnie z lematem~\ref{lem:os-obrotu}. Przeliczalna unia zbiorów dwuelementowych jest co najwyżej przeliczalna: wyliczamy elementy grupy i dla każdego wypisujemy oba punkty osi. Można najpierw wypisać punkt, którego pierwsza niezerowa współrzędna jest dodatnia, więc ten porządek nie wymaga nowych arbitralnych wyborów.

Jeśli $gx=x$, to dla dowolnego $h\in G$ mamy $(hgh^{-1})(hx)=hx$. Sprzężony element $hgh^{-1}$ nie jest neutralny, gdyż równość $hgh^{-1}=e$ implikowałaby $g=e$. Zatem $hD\subseteq D$. Stosując to samo do $h^{-1}$, dostajemy $D\subseteq hD$, więc $hD=D$. Ponieważ $h$ jest bijekcją sfery, także $hX=X$.

Wreszcie, jeśli $x\in X$ i $gx=x$, to $g\ne e$ oznaczałoby $x\in D$, sprzecznie z definicją $X$. Stabilizator każdego punktu $X$ jest więc trywialny.
\end{proof}

\subsection{Wybór reprezentantów: dokładne miejsce aksjomatu wyboru}
\begin{definicja}[Aksjomat wyboru w potrzebnej postaci]
Dla każdej rodziny niepustych zbiorów istnieje funkcja, która każdemu zbiorowi z tej rodziny przyporządkowuje jeden z jego elementów.
\end{definicja}
Zastosujmy ten aksjomat do rodziny wszystkich orbit działania $G$ na $X$. Wybieramy zbiór $T\subseteq X$ zawierający dokładnie po jednym punkcie z każdej orbity. To jest zasadniczy niekonstruktywny wybór w dowodzie. Nie podajemy algorytmu wyznaczającego $T$.
\begin{lemat}[Jednoznaczne współrzędne orbitowe]\label{lem:wspolrzedne-orbitowe}
Każdy punkt $x\in X$ ma dokładnie jedno przedstawienie $x=gt$, gdzie $g\in G$ i $t\in T$.
\end{lemat}
\begin{proof}
Istnienie wynika z wyboru reprezentanta orbity $x$. Jeśli $gt=hu$, to $t$ i $u$ należą do tej samej orbity, bo $t=g^{-1}hu$. Z wyboru jednego reprezentanta wynika $t=u$. Równość $gt=ht$ daje $(h^{-1}g)t=t$. Swoboda działania wymusza $h^{-1}g=e$, czyli $g=h$.
\end{proof}
\begin{uwaga}
To, że każda orbita jest przeliczalna, nie daje automatycznie przepisu wyboru jednego punktu z każdej orbity. Orbit może być nieprzeliczalnie wiele. Aksjomat wyboru nie jest tu skrótem dla wyliczenia wszystkich punktów sfery.
\end{uwaga}

\subsection{Cztery kawałki na wszystkich orbitach jednocześnie}
Dla $C\subseteq G$ definiujemy $C^T=\{gt:g\in C,\ t\in T\}\subseteq X$. Z jednoznaczności w lemacie~\ref{lem:wspolrzedne-orbitowe} wynika, że rozłączne podzbiory $G$ dają rozłączne zbiory $C^T$. Ponadto
\begin{equation}\label{eq:6-2}
 h(C^T)=(hC)^T\qquad(h\in G),
\end{equation}
ponieważ $h(gt)=(hg)t$. Na wszystkich orbitach działa ten sam obrót $h$, niezależny od wybranego reprezentanta.
\begin{twierdzenie}[Paradoksalność $X$]\label{tw:paradoksalnosc-x}
Zbiór $X=S^2\setminus D$ jest paradoksalny pod działaniem obrotów. Wystarczają cztery części.
\end{twierdzenie}
\begin{proof}
Przenosimy zbiory z \eqref{eq:4-3} na orbity: $C_1=A_1^T$, $C_2=A_2^T$, $C_3=B_1^T$, $C_4=B_2^T$. Tworzą partycję $X$, ponieważ odpowiadające im zbiory grupowe tworzą partycję $G$, a każdy punkt ma jednoznaczne współrzędne $gt$. Z \eqref{eq:4-4} i \eqref{eq:6-2} otrzymujemy $X=C_1\sqcup aC_2$ oraz $X=C_3\sqcup bC_4$. To rozkład $X\equivG2X$. Ruchy to $h_1=I$, $h_2=A$, $h_3=I$, $h_4=B$, a etykiety docelowe to odpowiednio $1,1,2,2$.
\end{proof}
Nie mnożymy liczby kawałków przez liczbę orbit: $C_1$ jest jednym zbiorem zawierającym stosowne punkty wszystkich orbit, i tak samo dla $C_2,C_3,C_4$.

